% ============================================================
%  ae Estimation MLP — Hyperparameter Optimisation Analysis
%  Thomas M. Rudolf
% ============================================================
\documentclass[11pt, a4paper]{article}

\usepackage[utf8]{inputenc}
\usepackage[T1]{fontenc}
\usepackage[english]{babel}
\usepackage{geometry}
\geometry{margin=2.5cm}

\usepackage{amsmath, amssymb}
\usepackage{booktabs}
\usepackage{graphicx}
\usepackage{hyperref}
\usepackage{cleveref}
\usepackage{natbib}
\usepackage{microtype}
\usepackage{setspace}
\usepackage{parskip}

\setstretch{1.15}

\hypersetup{
  colorlinks = true,
  linkcolor  = blue!60!black,
  citecolor  = blue!60!black,
  urlcolor   = blue!60!black
}

% ---- convenience macros ----
\newcommand{\Lhuber}{\mathcal{L}_{\text{Huber}}}
\newcommand{\ae}{a_e}
\newcommand{\MAE}{\text{MAE}}
\newcommand{\Rtwo}{R^2}

% ============================================================
\begin{document}

\title{\textbf{Hyperparameter Optimisation Results}\\[4pt]
       \large \textit{ae} Estimation MLP: Grid Search vs.\ Bayesian TPE}
\author{Thomas M. Rudolf}
\date{\today}
\maketitle
\tableofcontents
\bigskip

% ============================================================
\section{Overview}
\label{sec:overview}

Three complementary diagnostic figures were generated to evaluate the effectiveness
of Bayesian Tree-structured Parzen Estimator (TPE) optimisation against a full
Grid Search for tuning the Multi-Layer Perceptron (MLP) used to estimate the radial
depth of cut ($\ae$) from spindle motor current signals.
The figures address (i) the aggregate performance and convergence of both strategies,
(ii) the sensitivity of validation loss to pairwise hyperparameter combinations
identified via Grid Search, and (iii) the functional variance decomposition (fANOVA)
of hyperparameter importance as computed by the TPE sampler.
Together, they provide a complete characterisation of the optimisation landscape
governing the MLP's generalisation behaviour.

% ============================================================
\section{Figure~1 — Grid Search vs.\ Bayesian TPE: Overall Performance}
\label{sec:fig1}

% ------ 1A ------
\subsection{Panel~(A) — Validation Loss Distributions}
\label{sec:fig1a}

The Grid Search explored an exhaustive grid of $n = 216$ configurations, whereas
Bayesian TPE required only $n = 80$ trials — a 63\% reduction in computational
expenditure.
The Grid Search distribution is tightly concentrated between
$\Lhuber \approx 0.025$ and $0.060$, indicating that a narrow effective region in
the hyperparameter space was identified consistently across configurations.
The TPE distribution is noticeably wider, extending past $0.150$, which reflects
the expected exploratory behaviour during early trials before the surrogate model
has accumulated sufficient observations to guide sampling reliably.

Critically, the dashed medians of both distributions converge to essentially the
same best-region performance ($\Lhuber \approx 0.022$), confirming that TPE
reaches the global optimum at a fraction of the Grid Search cost.
The heavier right tail of the TPE distribution is not a concern: it is a structural
artefact of random initialisation in trials 1–20, before surrogate guidance becomes
effective.

% ------ 1B ------
\subsection{Panel~(B) — TPE Convergence Trajectory}
\label{sec:fig1b}

The ``best so far'' curve drops steeply within the first 15 trials, indicating
that the surrogate model locates the high-performing region of the hyperparameter
space very early in the search.
By trial $\approx 55$, the best-so-far value effectively reaches the Grid Search
best of $\Lhuber = 0.0221$ (shown as the blue dashed reference line).
No further improvement is observed beyond trial~55, confirming that the optimisation
has saturated.
This is a healthy convergence profile: steep early gains followed by a stable
plateau.
The implication is that 55--60 TPE trials would have been sufficient; the final
20 trials served only to confirm convergence rather than to discover new optima.

% ------ 1C/D ------
\subsection{Panels~(C) and~(D) — Parity Plots}
\label{sec:fig1cd}

Both the Grid Search best configuration and the TPE best configuration yield
identical $\Rtwo = 0.984$, confirming that the two optimal architectures are
functionally equivalent from a predictive standpoint.
However, TPE achieves a mean absolute error of $\MAE = 0.377\,\text{mm}$ compared
to $\MAE = 0.453\,\text{mm}$ for the Grid Search optimum — a 17\% reduction in
absolute estimation error.

Both models exhibit the greatest scatter at $\ae = 40\,\text{mm}$ (the largest
radial engagement level tested), where prediction variance increases.
This is physically consistent: higher $\ae$ values involve more complex engagement
dynamics, and a larger fraction of the cutting torque signal approaches the motor's
rated load.
At $\ae \approx 26\,\text{mm}$, both models exhibit a slight underprediction
tendency, which may reflect a reduced signal-to-noise ratio at low radial engagement.
The systematic point clustering visible at $\ae \in \{26, 32, 38, 40\}\,\text{mm}$
confirms that the experimental design employs discrete $\ae$ levels; no
interpolation across unseen engagement values is assessed in the current validation
set.

% ============================================================
\section{Figure~2 — Grid Search Sensitivity Heatmaps}
\label{sec:fig2}

% ------ 2 left ------
\subsection{Left Panel — Dropout $\times$ n\_layers}
\label{sec:fig2left}

The global minimum ($\Lhuber = 0.022$) occurs at
$\text{dropout} = 0.30$, $\text{n\_layers} = 3$.
This is a strong, isolated optimum: no other combination in the heatmap approaches
this value.
Performance degrades substantially toward $\text{dropout} = 0.45$ with
$\text{n\_layers} = 1$ ($\Lhuber = 0.028$) or $\text{n\_layers} = 2$
($\Lhuber = 0.027$), indicating that shallow networks with high dropout are
over-regularised and lose the representational capacity required for
physics-informed feature regression.

Examining the column-wise trend, $\text{n\_layers} = 3$ consistently yields the
lowest loss across all dropout levels, confirming that depth is the dominant
architectural factor in this panel.
Moderate dropout ($0.15$–$0.30$) acts as a beneficial regulariser, while
$\text{dropout} = 0.45$ degrades performance regardless of depth.
It has to be mentioned that the absence of an $\text{n\_layers} = 4$ condition
in the Grid Search represents a search boundary; whether additional depth would
yield further improvement remains an open question.

% ------ 2 right ------
\subsection{Right Panel — Units per Layer $\times$ n\_segments}
\label{sec:fig2right}

The minimum ($\Lhuber = 0.022$) appears at
$\text{units} = 64$, $\text{n\_segments} = 2$.
The most important observation is that performance degrades monotonically as
$\text{n\_segments}$ increases beyond~2, regardless of units per layer.
At $\text{n\_segments} = 7$, all configurations exceed $\Lhuber = 0.048$ —
more than double the optimal loss.
This reveals a clear over-segmentation effect: excessive temporal fragmentation
of the current signal destroys the structural information that the MLP relies upon
to reconstruct $\ae$.

The optimum at $\text{n\_segments} = 2$ is physically interpretable: two segments
plausibly correspond to the engagement phase and the steady-state cutting phase of
a tooth pass, a decomposition that carries genuine mechanical meaning.
Units per layer $= 128$ offers only marginal improvement over units $= 64$ at
$\text{n\_segments} = 2$ ($\Lhuber = 0.025$ vs.\ $0.022$), while units $= 32$
consistently underperforms across all segment counts.
This confirms that network capacity matters, but only up to a threshold of
64 units per layer, beyond which returns diminish rapidly.

% ============================================================
\section{Figure~3 — TPE fANOVA Hyperparameter Importance}
\label{sec:fig3}

Dropout dominates the fANOVA decomposition with an importance score of
approximately $0.55$ — exceeding the combined contribution of all remaining
hyperparameters.
This is the central finding of the optimisation: regularisation strength is the
primary determinant of generalisation performance in this MLP, not architectural
size or temporal resolution strategy.
n\_segments and n\_layers share roughly equal importance ($\approx 0.19$ each),
confirming that both the temporal feature extraction approach and the network
depth exert comparable influence on validation loss.
units\_l0 contributes only $\approx 0.06$, consistent with the heatmap observation
that model performance is relatively insensitive to layer width once the minimum
capacity threshold of 64 units is satisfied.

The fANOVA result also implies that future ablation studies or continued
optimisation should prioritise dropout range exploration before investing in
architectural search.
Specifically, the near-equivalence of n\_segments and n\_layers importance suggests
that both could be jointly tuned in a lower-dimensional follow-up search without
loss of coverage.

% ============================================================
\section{Synthesis}
\label{sec:synthesis}

The three figures together articulate a coherent optimisation narrative.
The MLP's generalisation capacity is governed primarily by dropout regularisation,
followed comparably by temporal segmentation strategy (n\_segments) and network
depth (n\_layers), with layer width contributing minimally beyond a capacity floor.
Bayesian TPE identified the optimal configuration —
$\text{dropout} = 0.30$, $\text{n\_layers} = 3$, $\text{units} = 64$,
$\text{n\_segments} = 2$ — using 63\% fewer evaluations than Grid Search, while
achieving a 17\% reduction in $\MAE$ ($0.377\,\text{mm}$ vs.\ $0.453\,\text{mm}$).
The parity plots confirm an excellent $\Rtwo = 0.984$ for both methods, but the
TPE-tuned model's lower $\MAE$ is the operationally relevant metric for
in-process $\ae$ estimation.

The optimum at $\text{n\_segments} = 2$ deserves particular emphasis: it carries a
physics-based interpretation tied to the tooth engagement phases of the milling
cycle, which strengthens the argument that the feature engineering strategy captures
genuine process information rather than spurious signal correlations.
This eliminates the need to rely solely on empirical justification for the temporal
segmentation choice and connects the data-driven model back to the underlying
cutting mechanics.

% ============================================================
\end{document}
